% ============================================================
%  Chapter 5 — Conclusion
%  Simple sentences throughout. Matches the actual notebook work.
% ============================================================

\chapter{Conclusion}
\label{ch:conclusion}

% --------------------------------------------------------
\section{Summary of Work}
\label{ch5-sec-summary}
% --------------------------------------------------------

This thesis built a feature-selection-based machine learning framework for predicting postpartum depression (PPD) risk in Bangladeshi mothers.
The central contribution is a principled pipeline that goes from a 70-column survey dataset to an interpretable, deployable prediction tool in five steps: systematic feature selection, a 15-model comparison, statistical risk factor analysis, SHAP explainability, and a web interface.

The dataset contains 800 Bangladeshi postpartum mothers and 70 columns, collected through a structured questionnaire survey.
After removing the serial number column and identifying six feature categories, the dataset has 69 usable columns divided into: Demographic (8 features), Familial (9 features), Personal Health (8 features), Neonatal Health (21 features), Screening Items (19 EPDS and PHQ-9 questions), and Target/Score (5 label columns).
The primary target variable is \textbf{EPDS Result}, a three-class categorical label: \textbf{Low} (EPDS Score 0--9), \textbf{Medium} (10--12), and \textbf{High} (13--30).

The work proceeded through seven main stages.
First, the raw data was cleaned through specific steps including fixing spelling inconsistencies, correcting entry errors, standardising category labels, and correcting 43 mislabelled EPDS Result entries.
Second, all 70 columns were mapped to one of six feature categories.
Third, all categorical columns were encoded using LabelEncoder.
Fourth, a Chi-square feature ranking and backward elimination experiment, using both Random Forest and Gradient Boosting as validation models, identified the \textbf{17 most predictive features} from 65 candidates --- all EPDS or PHQ-9 questionnaire items.
Fifth, all 15 machine learning and deep learning models were trained on these 17 features to compare their performance, with Gradient Boosting achieving the highest accuracy of 0.926.
Sixth, each of the four background feature categories was tested separately to measure its predictive power: Neonatal Health (61\%), Personal Health (52\%), Familial (47\%), and Demographic (40\%).
Seventh, SHAP analysis using a Permutation Explainer identified the top questionnaire items driving predictions for each EPDS risk class.
Eighth, Pearson, Spearman, and Kendall correlation coefficients were computed for the 17 selected features against the EPDS Score.
Ninth, within-category feature importance was computed by combining Random Forest, XGBoost, Logistic Regression, and Mutual Information scores.
Tenth, Kendall's tau was computed for all background features to identify the strongest statistical risk factors for PPD.

% --------------------------------------------------------
\section{Answers to the Research Objectives}
\label{ch5-sec-objectives}
% --------------------------------------------------------

Chapter~\ref{ch:introduction} stated five research objectives.
This section reports how each one was answered.

\textbf{Objective 1: Feature selection for PPD classification.}

Chi-square ranking combined with backward elimination identified 17 optimal EPDS and PHQ-9 questionnaire items as the most predictive subset from 65 candidate features.
These 17 items --- 10 EPDS questions and 7 PHQ-9 questions --- give 0.926 accuracy on the three-class EPDS Result target, which is higher than using all 19 questionnaire items without selection.
The four background categories, tested separately for comparison, achieve 40--61\% accuracy, far below the selected features.
This answers Objective~1 clearly: systematic feature selection improves over indiscriminate use of all items, and the selected EPDS and PHQ-9 items are far more predictive than background information alone.

\textbf{Objective 2: Machine learning vs.\ deep learning comparison.}

Traditional machine learning ensemble methods consistently outperform deep learning models.
Gradient Boosting achieves the highest accuracy (0.926) and the best performance in every experiment.
LightGBM, XGBoost, and Hist Gradient Boosting follow closely in the top four, all achieving accuracy above 0.916.
The three deep learning models (1D CNN, MLP, Deep MLP) rank 8th, 11th, and 15th respectively.
The performance gap between the best ensemble model and the best deep learning model is approximately 2.6\% in accuracy.
This result is consistent with other small-sample survey studies: with $n = 800$ and mostly categorical features, tree-based ensembles are a better structural fit than neural networks.
This answers Objective~2 clearly: deep learning does not offer a practical advantage over traditional ML for this Bangladeshi PPD survey dataset.

\textbf{Objective 3: Risk factor analysis and feature importance.}

From the Kendall's tau analysis of background features, the strongest risk factors are:
\begin{enumerate}
    \item \textbf{Anger after latest child birth} (Familial, $\tau = +0.418$, $p = 1.4 \times 10^{-45}$) --- the single strongest background predictor.
    \item \textbf{Feeling about motherhood} (Familial, $\tau = +0.280$, $p = 5.3 \times 10^{-22}$) --- negative feelings about becoming a mother are a strong risk signal.
    \item \textbf{Abuse} (Familial, $\tau = -0.223$, $p = 1.1 \times 10^{-14}$) --- the negative direction reflects under-reporting bias.
    \item \textbf{Depression during pregnancy PHQ2} (Neonatal Health, $\tau = +0.195$, $p = 3.3 \times 10^{-11}$) --- prenatal depression strongly predicts postnatal depression.
    \item \textbf{Trust and share feelings} (Familial, $\tau = -0.177$, $p = 1.8 \times 10^{-9}$) --- a protective factor.
\end{enumerate}

From the within-category importance analysis, the top feature inside each background category is: ``Anger after latest child birth'' in Familial, ``Depression during pregnancy (PHQ2)'' in Neonatal Health, ``History of depression'' in Personal Health, and ``Age'' in Demographic.
The dominant finding is that \textbf{emotional and familial factors matter far more than demographic factors} for background-based PPD risk.

\textbf{Objective 4: SHAP explainability analysis.}

SHAP analysis using a Permutation Explainer on the trained Gradient Boosting model identified the following as the three most important questionnaire items by mean absolute SHAP value:
\begin{enumerate}
    \item ``You have been so unhappy that you have been crying'' (EPDS, SHAP = 0.0585)
    \item ``You have felt scared or panicky for no good reason'' (EPDS, SHAP = 0.0500)
    \item ``You have felt sad or miserable'' (EPDS, SHAP = 0.0457)
\end{enumerate}
EPDS items overall outrank PHQ-9 items in importance.
Class-specific beeswarm plots show that high values on emotional symptom items push strongly toward the High EPDS risk class, while high values on the laughter item push toward the Low class.
This answers Objective~4: SHAP analysis provides an item-level, class-specific explanation of the model's predictions.

\textbf{Objective 5: Web interface for practical deployment.}

A web interface was built and deployed using the saved Gradient Boosting model (\texttt{gb\_model.pkl}) and scaler (\texttt{gb\_scaler.pkl}).
The interface accepts responses to the 17 selected EPDS and PHQ-9 questionnaire items and outputs an EPDS Result classification (Low, Medium, or High) with a brief clinical interpretation.
This answers Objective~5: the trained model is accessible to healthcare workers without requiring any programming knowledge.

% --------------------------------------------------------
\section{Contributions}
\label{ch5-sec-contributions}
% --------------------------------------------------------

This thesis makes five contributions to the PPD prediction literature.

\textbf{Contribution 1: First systematic feature selection study on this 800-mother Bangladeshi PPD dataset.}
Previous studies on Bangladeshi PPD data used all questionnaire items as input without selection.
This thesis applies Chi-square ranking combined with backward elimination to identify the 17 most predictive items, reducing the questionnaire burden from 19 to 17 items and improving accuracy over the unselected baseline.

\textbf{Contribution 2: Systematic comparison of 15 ML and deep learning models.}
No prior Bangladeshi PPD study compared 15 models including both traditional ML and three different deep learning architectures on the same dataset under the same cross-validation conditions.
This comparison produces a clear answer: Gradient Boosting is the best model, and deep learning does not outperform traditional tree ensembles for this type of data.

\textbf{Contribution 3: Statistically grounded risk factor analysis using Kendall's tau.}
This thesis uses Kendall's tau, a rank-based statistical test, to measure the association between each background feature and the EPDS Score independently of any specific model.
This gives a model-agnostic picture of which background risk factors matter most in the Bangladeshi population.
Anger after childbirth ($\tau = +0.418$) emerges as the single strongest background predictor, a finding that is robust across all methods used in this thesis.

\textbf{Contribution 4: SHAP explainability at the questionnaire-item level.}
This thesis applies a SHAP Permutation Explainer to the trained Gradient Boosting model to identify which of the 17 selected questionnaire items drive predictions for each EPDS risk class.
The top three items (crying, panic, sadness) provide an immediately actionable triage signal.
The item-level SHAP explanation makes the model interpretable for clinical use.

\textbf{Contribution 5: Web interface for practical deployment.}
The trained Gradient Boosting model is deployed as a web interface that accepts responses to the 17 selected questionnaire items and outputs an EPDS risk classification.
This is the first deployable PPD prediction tool built from this 800-mother Bangladeshi dataset.

% --------------------------------------------------------
\section{Limitations}
\label{ch5-sec-limitations}
% --------------------------------------------------------

This thesis has four limitations.

\textbf{Limitation 1: Sample size and geographic coverage.}
The dataset contains 800 mothers, likely collected from a limited number of healthcare sites.
Results may not generalise to all regions of Bangladesh, particularly rural areas where healthcare access, social norms, and economic conditions differ significantly.

\textbf{Limitation 2: Cross-sectional design.}
All data was collected at one point in time.
The study can confirm associations but cannot establish causal direction.
A longitudinal study following mothers over time would address this limitation.

\textbf{Limitation 3: Self-reported data and social desirability bias.}
All features are self-reported through a structured questionnaire.
The abuse direction paradox discussed in Chapter~\ref{ch:results} shows how self-report can diverge from clinical reality when the question covers a stigmatised topic.

\textbf{Limitation 4: SHAP computed on a sample.}
The SHAP Permutation Explainer was run on a sample of 300 instances for computational feasibility.
Running on the full 800-mother dataset would give a more precise feature ranking, though the top features are unlikely to change substantially.

% --------------------------------------------------------
\section{Future Work}
\label{ch5-sec-future}
% --------------------------------------------------------

Several directions are open for future work.

\textbf{Multi-site data collection.}
Recruiting mothers from multiple hospitals and community health centres across different divisions of Bangladesh would improve the diversity and representativeness of the dataset and allow more robust model evaluation.

\textbf{Longitudinal follow-up study.}
Collecting data from the same mothers at multiple time points would allow the study of how PPD risk changes over time and whether early risk factor measurements can predict later PPD onset more accurately.

\textbf{Shortened screener from SHAP findings.}
The SHAP analysis suggests that the top 5--7 EPDS items carry most of the predictive power.
A future study could test whether a shortened screener using only the top SHAP-ranked items achieves comparable accuracy to the full 17-item version, reducing the clinical burden further.

\textbf{Multi-class SHAP deployment.}
The web interface currently outputs the overall EPDS class prediction.
A future version could display the class-specific SHAP waterfall chart for each individual prediction, helping clinicians understand which items drove the classification for that specific mother.

\textbf{Bengali-language mobile screening tool.}
The findings of this thesis, especially the five key background risk factors identified by Kendall's tau and the 17-item screener from feature selection, could power a Bengali-language mobile application for community health workers in Bangladesh.

\textbf{Validation in similar South Asian populations.}
Bangladesh shares many social, cultural, and economic characteristics with India, Pakistan, and Nepal.
Testing whether the risk factor findings and feature selection results replicate in those populations would help assess the generalisability of the results.

% --------------------------------------------------------
\section{Closing Remarks}
\label{ch5-sec-closing}
% --------------------------------------------------------

Postpartum depression affects nearly half of all new mothers in some Bangladeshi healthcare settings, yet most cases go undetected because of stigma, lack of awareness, and a shortage of trained mental health professionals.
Machine learning can help bridge this gap --- but only when it is applied carefully, with principled feature selection and meaningful explainability.

This thesis demonstrates that systematic Chi-square feature selection combined with backward elimination can identify a compact, high-accuracy screener from EPDS and PHQ-9 questionnaire data.
The 17 selected items achieve 0.926 accuracy for three-class EPDS Result classification.
The SHAP analysis reveals that crying, fear or panic, and sadness are the three most important items --- a finding that is immediately actionable for clinicians.
The Kendall's tau analysis further reveals that anger after childbirth is the single strongest background predictor of PPD, a simple observable signal that a community health worker can ask about without any formal questionnaire.

A mother who reports strong anger, negative feelings about motherhood, and lack of a trusted support person in her family is at high statistical risk of developing PPD and should be referred for early professional assessment.
The web interface built in this thesis makes the trained model accessible for exactly that purpose.
